\section{Regularity and continuation at the transmission surface}
\label{sec:transmission-analysis}

We prove the existence assertion needed in Section~\ref{sec:compact-construction}.
Throughout this section all parameters, including the penalty strengths, are
finite.  Constants in estimates for the lapse and height state may depend on
the fixed parameters preceding the strengths, but not on those strengths.
Constants involving the conformal factor or higher derivatives may also
depend on the strengths.  Neither assertion requires uniform ellipticity as
the internal parameters tend to their final limits.

Write $\Omega_R=\Omega_i\cup_{\mathcal C}\Omega_{o,R}$, with far boundary
$S_R$.  On both sides of $\mathcal C$ the normal $n$ points toward the end.
Thus equality of the two $n$-fluxes is the condition for cancellation in
Green's formula.  The two base metrics have the same induced metric at
$\mathcal C$.  Joined Fermi charts have coordinates $(z',t)$, with
$t<0$ inside, $t>0$ outside, and $n=\partial_t$ at the face.  Coordinate
divergence currents below include the relevant volume density.  We retain
the notation $d=s/N$, $y=1-d^2$, $q=(1-y)/(1+y)$,
$I=(1+y)^{-1}$, $J=y/(1+y)$, and $\beta=2/(1+y)$ from the graph calculus.

\begin{proposition}[Finite-parameter transmission existence]
\label{prop:transmission-existence}
Fix the smooth compact base data and AF attachment of
Section~\ref{sec:compact-construction}, with $H_i>|k_T|$ on
$\mathcal C$, $H_g(S)\leq0$, and the smooth negative definite synthetic
tensor $K$.  Fix positive finite parameters and smooth cutoffs as in
\eqref{eq:compact-system}--\eqref{eq:compact-trace-cutoff}, including a
source-free exterior collar of $\mathcal C$ and the constant $b_*>0$.
Then on every sufficiently large smooth truncation there is a positive
piecewise smooth solution of \eqref{eq:compact-system} and
\eqref{eq:compact-boundary} at $u_*=1$.

At these fixed parameters a subsequence of such solutions as $R\to\infty$
converges smoothly on each closed side on compact subsets to a positive
piecewise smooth solution on the full joined exterior.  Its end values are
$N_o,w\to1$.  The bounds in \eqref{eq:compact-bounds} hold, and all
compact-set derivative bounds are uniform in $R$.  The penalty conclusions
in Section~\ref{sec:compact-construction}, when the strengths have been
chosen there, therefore hold for these solutions.
\end{proposition}

The proof uses the a priori estimates of Lemma~\ref{lem:compact-lapse},
which were proved for positive classical zeros without assuming existence.
The scalar linear transmission fact used in that lemma is established
independently below as part of Lemma~\ref{lem:transmission-fredholm}.
In particular there is no appeal to
Proposition~\ref{prop:compact-solutions} in proving this proposition.

\subsection{Energy and oscillation of the lapse and height state}

For the moment consider any positive homotopy zero in the continuation
domain: $f\in C^{3,\alpha'}$ on the closed inner side and
$N,w\in C^{2,\alpha'}$ on each closed side, with $\alpha'>0$.
No bound for these norms is assumed.  The exponent $\gamma$ obtained
below depends only on the finite state bounds, so the continuation
domain can subsequently be fixed with $\alpha'<\gamma$.
The finite bounds in \eqref{eq:compact-bounds} give positive constants
$c_0,C_0$ such that
\begin{equation}
 c_0\leq N_i,s_i,N_o\leq C_0,
 \qquad |\dd f|\leq C_0.
 \label{eq:transmission-state-range}
\end{equation}
Here and below an estimate for $s$ on the exterior means the same estimate
for $N_o$.  The scalar lapse sources and the drift
$(u_*K+\ell g)p$ are bounded at this stage; their derivatives are not used.

\begin{lemma}[Continuity at the nonlinear seam]
\label{lem:seam-holder}
For the state range \eqref{eq:transmission-state-range}, there are
$\gamma>0$ and $C<\infty$ such that
\begin{equation}
 \|N_i\|_{C^\gamma(\overline\Omega_i)}
 +\|\dd f\|_{C^\gamma(\overline\Omega_i)}
 +\|N_o\|_{C^\gamma(K\cap\overline\Omega_o)}\leq C
 \label{eq:transmission-holder}
\end{equation}
on every fixed compact set $K$.  The constants are uniform over homotopy
zeros and sufficiently large outer truncations.  They depend on the finite
state range, bounds for the undifferentiated height and lapse sources,
and the fixed geometry, but not on derivatives of $T$ or the strength of
$P$.
\end{lemma}

\begin{proof}
Away from $\mathcal C$, the lapse equation is a smooth metric Laplace
equation with bounded divergence forcing and bounded scalar forcing.
Its local $W^{1,p}$ bounds hold for some $p>3$, also up to $S$ with its
natural current flux.  The prescribed constant flux at $S$ can be put in
the current by a smooth bounded extension before reflection.  Differentiating
the divergence height equation gives scalar uniformly elliptic equations
for the coordinate derivatives of $f$.  Their divergence forcing contains
first derivatives of $N$, bounded in $L^p$, and the bounded scalar height
forcing, put into the corresponding component of the differentiated
current.  In particular differentiating that scalar forcing in the weak
equation does not require a bound for $\dd T$.
The inhomogeneous scalar divergence estimate gives H\"older bounds for
$\dd f$.  At $S$ subtract the constant height and reflect it oddly, the
lapse evenly, and the metric with its tensor parities in an orthogonal
collar.  The reflected metric is Lipschitz.  The height flux satisfies the
weak reflected equation, since its normal component has the required even
parity.  Difference quotients justify differentiation there.  These
arguments prove the assertion away from the seam.

At the seam we first bound energy weighted by inverse distance.
A modified Newton current then gives an estimate relative to compatible
constant states. The weighted bound supplies a uniformly positive
scale of small excess; a blowup to an explicit linear transmission
problem gives excess decay and hence H\"older continuity.

\smallskip\noindent\emph{Weighted energy.}
Set
\begin{equation}
 E=\begin{cases}|\nabla N_i|^2+|\nabla^2 f|^2,&t<0,\\
                 |\nabla N_o|^2,&t>0,
   \end{cases}
 \qquad G=-\log s,\qquad Y=s^{-k_0},
 \label{eq:transmission-energy-definition}
\end{equation}
where $k_0$ is sufficiently large for
Lemma~\ref{lem:gradient-energy}.  On the inside its divergence matrix is
$\rho\mathcal P$, where $\rho=(1+y)/d$, and on the outside it is the
base cometric.  After including the volume densities, these give a bounded
symmetric uniformly elliptic matrix $a$ in the joined chart.
$Y$ is continuous across the face because $N_o=s_i$ there.
The energy identity and the boundary identity give
\begin{equation}
 \partial_j(a^{jk}\partial_kY)
 \geq cE-C-C\delta_{\mathcal C}.
 \label{eq:transmission-joined-energy}
\end{equation}
This is a distributional inequality.  Indeed $\rho q=d$ on a
Dirichlet-height face, and the inner normal $Y$-flux is
$-k_0YB_n/s+O(1)$, whereas the outer one is $-k_0YB_n/s$.
The bounded discrepancy is precisely the surface term in
\eqref{eq:transmission-joined-energy}.

Fix concentric chart balls strictly contained in a larger joined chart.
Let $\theta_\varepsilon$ be the zero Dirichlet Green potential, for
$-\partial(a\partial)$ on the larger ball, of a nonnegative normalized
bump at a point $z_0$ in the smaller ball.  If $Y_H$ is the homogeneous
replacement of $Y$ with its boundary values, then
\begin{equation}
 \big\langle\partial(a\partial Y),\theta_\varepsilon\big\rangle
  =\int(Y_H-Y)\,\chi_\varepsilon\leq C.
 \label{eq:transmission-green-pairing}
\end{equation}
The maximum principle bounds $Y_H$, and $Y$ is bounded by
\eqref{eq:transmission-state-range}.  For each classical solution the
identity follows first with regular potentials and then by Sobolev
approximation.  The surface term is well defined by the trace theorem.
The divergence Green bounds for symmetric measurable uniformly elliptic
matrices in dimension three give an upper bound $C|z-z_0|^{-1}$ and a
comparable lower bound on a sufficiently small interior ball
\citep[Theorem~7.1 and pp.~66--67]{LSW1963}.  The surface integral of
$|z-z_0|^{-1}$ is uniformly finite.  Letting the bump shrink, using
\eqref{eq:transmission-joined-energy} and Fatou's lemma, proves
\begin{equation}
 \int_{B_{h_0}(z_0)}\frac{E(z)}{|z-z_0|}\,\dd z\leq C.
 \label{eq:transmission-green-energy}
\end{equation}
Both $h_0>0$ and $C$ are uniform on the finite collection of seam charts.

\smallskip\noindent\emph{Comparison with a constant state.}
We next prove an energy bound relative to a compatible constant state.
Such a state is
\begin{equation}
 \mathcal V_*=(n_*,a_*\dd t,s_*),\qquad
 s_*=s(n_*,|a_*|),\qquad
 \mathcal V=(N_i,\dd f,N_o).
 \label{eq:transmission-reference}
\end{equation}
The pairs $(n_*,a_*)$ range over a compact rectangle containing the
state ranges with a margin, with $n_*>0$.  Since $|\dd t|_g=1$, $s_*$ is
constant in this Fermi chart.  Norms of the state mean the sum of the
inner pair's squared norm on the inner half and the outer lapse's
squared norm on the outer half.

To derive the bound, put
$\mathcal T_f=\nabla^2f-(\Delta f)g$.  The Newton-current identity in
Lemma~\ref{lem:newton-flux} says
\begin{equation}
 \rho\mathcal P\nabla G
 =-\frac{\nabla N_i}{s}
   +\mathcal T_f\left(\frac{y\nabla f}{d|\dd f|^2}\right)+O(1).
 \label{eq:transmission-newton-current}
\end{equation}
The coefficient vector extends smoothly through $\dd f=0$.  In the
current $a\dd Y$, add $k_0Y_*B/s_*$ on both sides, where
$Y_*=s_*^{-k_0}$, and on the inside subtract the Newton term in
\eqref{eq:transmission-newton-current} with its full coefficient
$k_0Y,y\nabla f/(d|\dd f|^2)$ frozen at the reference.  Include metric
densities in all three currents.  Call the resulting joined current
$\mathcal J_*$.  The cancellation of the coefficients of $\nabla N$
and $\nabla^2f$ at the reference gives
\begin{equation}
 |\mathcal J_*|\leq C\bigl(1+|\mathcal V-\mathcal V_*|\sqrt E\bigr).
 \label{eq:transmission-modified-current-size}
\end{equation}
Moreover $\divg\mathcal T_f=\Ric(\nabla f,\cdot)$, so the divergence
of the subtracted term is bounded by $C(1+\sqrt E)$; its frozen vector
is extended smoothly using $\partial_t$.  The added $B$-current has
bounded divergence and no seam jump.  The frozen Newton vector is normal
at the seam.  Its normal contraction there involves
$-\tr_{\mathcal C}\nabla^2f$, which is bounded since $f$ is constant
along that face and $|\dd f|$ is bounded.  Young's inequality therefore
gives, after lowering $c$,
\begin{equation}
 \divg\mathcal J_*\geq cE-C-C\delta_{\mathcal C}.
 \label{eq:transmission-modified-current-divergence}
\end{equation}
Test by $\chi^2$, with $\chi=1$ on $B_h$, supported in $B_{2h}$ and
$|\dd\chi|\leq C/h$.  Absorb half the energy using
\eqref{eq:transmission-modified-current-size}.  The volume, constant
current, and surface contributions are at most $Ch^3$, $Ch^2$, and
$Ch^2$, respectively.  We obtain
\begin{equation}
 h^{-1}\int_{B_h}E\leq
 C\left(\frac{1}{|B_{2h}|}\int_{B_{2h}}
                    |\mathcal V-\mathcal V_*|^2+h\right).
 \label{eq:transmission-caccioppoli}
\end{equation}
On a pure inner ball the identical construction uses any constant
covector reference.  The coefficients frozen against this reference may
vary smoothly with the background metric; their derivatives give only
$C(1+\sqrt E)$ errors.  On a pure outer ball the analogous assertion
is the usual lapse Caccioppoli inequality.

\smallskip\noindent\emph{A scale of small excess.}
Fix $0<b_0<1$ and define the seam excess
\begin{equation}
 \mathcal E(z_0,h)=
 \inf_{\mathcal V_*}\frac{1}{|B_h|}\int_{B_h(z_0)}
          |\mathcal V-\mathcal V_*|^2+h^{b_0}.
 \label{eq:transmission-excess}
\end{equation}
There is a uniformly positive scale at which this excess is arbitrarily
small.  To see the quantitative content, put $h_j=2^{-j}h_0$.  For each
$z\ne z_0$ the sum of $h_j^{-1}$ over the indices with $|z-z_0|<h_j$
is at most $C|z-z_0|^{-1}$.  Thus
\eqref{eq:transmission-green-energy} bounds
$\sum_j h_j^{-1}\int_{B_{h_j}}E$ uniformly.  Given $\eta>0$, first
make the starting ceiling small enough that $h_0^{b_0}+h_0^2<\eta$.
Among a fixed finite number of its dyadic descendants there is a scale
with $h^{-1}\int_{B_h}E<\eta$.  The number depends only on the
Green bound and $\eta$, so the chosen scale has a uniform positive
lower bound.

Here is why this is also a small excess scale.  Scaled Poincar\'e and
the half-ball trace inequality bound both the volume deviations and the
mean-square face deviations from each side mean by
\begin{equation}
 C h^{-1}\int_{B_h}E+Ch^2.
 \label{eq:transmission-trace-poincare}
\end{equation}
The additional $Ch^2$ allows for Christoffel terms when replacing
covariant Hessians by coordinate derivatives of $\dd f$.  The
tangential trace of $\dd f$ is zero; hence the tangential components of
its mean are small.  The trace relation $N_o=s(N_i,\partial_t f)$ and
the uniform Lipschitz bound for $s$ on the state range show that the
outer mean differs from the value obtained from the two inner means by
at most the square root of \eqref{eq:transmission-trace-poincare}.
This provides a compatible reference with excess $C\eta$.  All
references obtained this way lie in a compact subrange with room for
the subsequent small updates.

\smallskip\noindent\emph{Excess decay.}
We claim that, for fixed sufficiently small $\vartheta>0$, whenever
$h$ and $\mathcal E(z_0,h)$ are sufficiently small,
\begin{equation}
 \mathcal E(z_0,\vartheta h)
 \leq 2\bigl(C\vartheta^2+\vartheta^{b_0}\bigr)
                    \mathcal E(z_0,h).
 \label{eq:transmission-excess-decay}
\end{equation}
The constant $C$ is independent of $\vartheta$.  We give the
compactness argument, including the value trace.

If the claim fails, choose violating balls with
$\sigma_j^2=\mathcal E(z_j,h_j)\to0$ and best references
$(n_j,a_j\dd t,s_j)$.  Translate and dilate these balls to unit size
using $z=z_j+h_j\zeta$, and rotate tangential coordinates to be
orthonormal at their centers.  Define the lapse differences divided by
$\sigma_j$, and the height difference by
\begin{equation}
 u_{i,j}=\frac{N_i-n_j}{\sigma_j},\qquad
 u_{o,j}=\frac{N_o-s_j}{\sigma_j},\qquad
 \phi_j(\zeta)=
       \frac{f(z_j+h_j\zeta)-a_jh_j\zeta_3}{h_j\sigma_j}.
 \label{eq:transmission-normalized-fields}
\end{equation}
The trace of $\phi_j$ is zero.  Their state has bounded $L^2$ norm on
the unit ball.  Applying \eqref{eq:transmission-caccioppoli} on smaller
balls gives locally bounded $H^1$ norms for $u_{i,j},u_{o,j}$ and
$\nabla\phi_j$.  Indeed its error after normalization is
$h_j/\sigma_j^2\leq h_j^{1-b_0}\to0$.
The $L^2$ norm of $\phi_j$ is controlled by its gradient and its zero
trace.  Rellich compactness and compactness of the $H^1$ trace into
$L^2$ on smaller face patches give strong local $L^2$ convergence of
the state in the bulk and on the face, and weak $H^1$ convergence.

Pass to a convergent reference $(n_*,a_*,s_*)$.  The derivative of
the height flux with respect to the gradient is $s_*^2\mathcal P_*$;
its lapse derivative is $s_*^2(\beta/n_*)a_*\partial_t$.
Taylor remainders divided by $\sigma_j$ have $L^1$ norm tending to
zero, since the unnormalized state difference is $\sigma_j$ times
an $L^2$-bounded function.  The same assertion holds on the face by
the trace bounds.  Smooth metric errors are $O(h_j/\sigma_j)$,
which tends to zero.  The height source and lapse drift enter weak
tests with this same factor; the scalar lapse source has the smaller
factor $O(h_j^2/\sigma_j)$.  Thus no derivative of a source is being
passed to the limit.  Joined tests have a common value on the face,
so lapse current matching passes to equality of the limiting normal
derivatives.  We arrive at
\begin{align}
 \Delta u_i&=\Delta u_o=0,
 & (q\partial_t^2+\Delta_{z'})\phi+k\partial_tu_i&=0,
 & k&=\beta a_*/n_*,\label{eq:transmission-frozen-interior}\\
 \phi&=0,
 &\partial_tu_i&=\partial_tu_o,
 &u_o&=r_0u_i+r_0c\partial_t\phi \quad(t=0),
 \label{eq:transmission-frozen-boundary}
\end{align}
where
\begin{equation}
 r_0=dI>0,\qquad c=-\frac{n_*y}{a_*},\qquad
 ck=q-1,\qquad r_0c=-\frac{s_*J}{a_*}.
 \label{eq:transmission-frozen-relations}
\end{equation}
All expressions have their continuous interpretations at $a_*=0$;
in particular $c=k=0$ there.  Differentiating
$s^2+|\dd f|^2s^4=N_i^2$ gives
$\dot s=dI\dot N_i-sJ\partial_t\dot f/a_*$, which verifies
the value relation directly.

The elementary harmonic estimate proved just below gives uniform
smooth bounds for these limits on smaller closed half-balls in terms
of their $L^2$ state norms.  Consequently their deviations from the
center values have mean square at most $C\vartheta^2$.  The center
values satisfy the linearized compatibility relation.  Replace the
old reference by
\begin{equation}
 n_j'=n_j+\sigma_j u_i(0),\qquad
 a_j'=a_j+\sigma_j\partial_t\phi(0),\qquad
 s_j'=s(n_j',|a_j'|).
 \label{eq:transmission-reference-update}
\end{equation}
The nonlinear compatibility error in the outer value is
$O(\sigma_j^2)$, and the tangential height gradient is zero at the
center.  Strong local $L^2$ convergence therefore proves
\eqref{eq:transmission-excess-decay}, contradicting its violation.
The added radius term contributes at most
$\vartheta^{b_0}\sigma_j^2$, as required.

For completeness, the harmonic estimate used here is as follows.
Reflect $u_o$ into $t<0$ and put $U=u_i+u_o^{\rm refl}$.  It is
harmonic with zero Neumann data, so it is smooth with controlled local
norms by even reflection.  On a smaller inner cylinder choose a
harmonic primitive $D$ with $D_t=u_i$.  An explicit choice is
\begin{equation}
 D(z',t)=\int_{t_0}^{t}u_i(z',v)\,\dd v+a(z'),
 \qquad \Delta_{z'}a=-\partial_tu_i(z',t_0),
 \label{eq:transmission-harmonic-primitive}
\end{equation}
where $t_0<0$ is a fixed interior slice.  The slice data are controlled
by interior harmonic estimates.  This formula gives $D\in H^2$
locally up to the face: the mixed and double normal derivatives come
from $u_i\in H^1$, and tangential Laplace estimates applied to
$\Delta_{z'}D=-\partial_tu_i$ give its remaining second derivatives.
The function $V=D+c\phi$ satisfies
$(q\partial_t^2+\Delta_{z'})V=0$, since $ck=q-1$, and $V=D$
on the face.  Set $\widetilde V(z',t)=V(z',\sqrt q\,t)$.
It is harmonic and, with $a_0=r_0/\sqrt q>0$,
\begin{equation}
 (\widetilde V-D)|_{t=0}=0,
 \qquad \partial_t(D+a_0\widetilde V)|_{t=0}=U.
 \label{eq:transmission-positive-harmonic-combinations}
\end{equation}
Scalar harmonic Dirichlet and Neumann estimates control both
combinations.  Their coefficient matrix is invertible because
$1+a_0>0$; in fact $a_0=1/\sqrt{1+y}$.  They control $D$, then
$u_i=D_t$ and $u_o=U-u_i^{\rm refl}$.  Finally the scalar Dirichlet
equation in \eqref{eq:transmission-frozen-interior} controls $\phi$.
This last step avoids division by $c$, so all estimates are uniform
through zero slope.  The normal rescaling only reduces the patch by
a uniform factor on the state range.

Choose $\vartheta$ so small, and then $\gamma>0$ so small, that
$2(C\vartheta^2+\vartheta^{b_0})\leq\vartheta^{2\gamma}$.
Begin iteration at the small-excess scales with uniform positive
lower bound constructed above.  The references stay in the permitted
range because their successive changes are bounded by a convergent
geometric sum.  This yields a mean-square oscillation bound
$Ch^{2\gamma}$ at every seam center.  On a pure ball near the seam,
the first scale comparable to its distance from the seam is controlled
by a seam ball.  The corresponding interior excess argument, with an
arbitrary constant gradient reference, then gives the bound at all
smaller scales.  Its limit equations are a harmonic lapse and a
constant uniformly elliptic height equation with the known lapse
coupling; scalar interior estimates give the same decay.  Campanato's
characterization on each closed half-chart proves
\eqref{eq:transmission-holder}.  Finitely many charts and the earlier
ordinary boundary estimates complete the proof.
\end{proof}

\subsection{Estimates with the full derivative orders}

We first record the linear estimate, including the unequal orders of
the unknowns.  This is also a proof of the relevant boundary estimate;
it does not invoke a general theorem for the nonlinear transmission
condition.

\begin{lemma}[Local transmission Schauder estimate]
\label{lem:transmission-schauder}
Give a height variation $\phi$ order $3$, lapse variations $u_i,u_o$
order $2$, and conformal-factor variations $v_i,v_o$ order $2$.
For frozen coefficients from a state in
\eqref{eq:transmission-state-range}, the principal equations are
\begin{equation}
 \mathcal P:\nabla^2\phi
       +\frac\beta N\nabla f\cdot\nabla u_i,\qquad
 \Delta u_i,\quad\Delta u_o,\qquad
 A_i:\nabla^2v_i,\quad A_o:\nabla^2v_o.
 \label{eq:transmission-principal-rows}
\end{equation}
Their residual norms are respectively $C^{1,\alpha}$ and
$C^{0,\alpha}$ for the other four rows.  The height trace has
$C^{3,\alpha}$ norm.  The lapse and conformal-factor value residuals
have $C^{2,\alpha}$ norm, and their current-flux residuals have
$C^{1,\alpha}$ norm.  On a smaller fixed patch, the field norms of
the indicated orders are bounded by these residual norms and the
supremum norms of the fields on a larger patch.  The constants are
uniform on compact positive state ranges.

At the seam the height/lapse principal conditions are
\eqref{eq:transmission-frozen-boundary}; the conformal-factor block
has value continuity and positive conormal matching.  At $S$ they
are Dirichlet for height and Neumann for the other fields; at $S_R$
they are Dirichlet.  The same estimates hold for scalar transmission
with a fixed positive $C^{2,\alpha}$ value ratio, and while its outer
normal-derivative coefficient is decreased to zero, the inner
coefficient being kept strictly positive.
\end{lemma}

\begin{proof}
It suffices first to work on a flat seam patch.  Subtract scalar
particular solutions of the two lapse Poisson equations, with
$C^{2,\alpha}$ control, and a $C^{3,\alpha}$ particular solution
of the height row, including its $C^{3,\alpha}$ boundary trace.
The height forcing includes the first derivative of the particular
inner lapse, hence belongs to $C^{1,\alpha}$ as required.  Such
particular solutions can be made on smooth domains having flat face
portions containing the smaller patch, after extending localized
data; ordinary scalar Dirichlet estimates give their stated norms.
We are left with harmonic lapses and the homogeneous height row,
zero height trace, and residuals
\begin{equation}
 \partial_tu_i-\partial_tu_o=h_1,
 \qquad u_o-r_0u_i-r_0c\partial_t\phi=h_2,
 \label{eq:transmission-inhomogeneous-seam}
\end{equation}
where $h_1\in C^{1,\alpha}$ and $h_2\in C^{2,\alpha}$.
The reflected sum $U=u_i+u_o^{\rm refl}$ is harmonic with Neumann
data $h_1$, and has a controlled $C^{2,\alpha}$ norm on a smaller
patch.  Construct $D$ by
\eqref{eq:transmission-harmonic-primitive} and use its same local
supremum bound.  The two harmonic combinations in
\eqref{eq:transmission-positive-harmonic-combinations} now have
Dirichlet datum zero and Neumann datum $U-h_2$, respectively.
The latter belongs to $C^{2,\alpha}$, so both combinations, and
therefore $D$, are controlled in $C^{3,\alpha}$.  This controls
both lapses in $C^{2,\alpha}$.  Scalar Dirichlet estimates for the
height row finish the $C^{3,\alpha}$ estimate for $\phi$.

For the conformal-factor block, the frozen inner matrix in orthonormal
face coordinates is a positive multiple of
$\operatorname{diag}(1,1,d^2)$, and the outer matrix is isotropic.
Subtract particular solutions, rescale the two normal coordinates,
and reflect the outer field.  The resulting fields are harmonic;
their difference has the prescribed value datum and a positive
weighted sum has the prescribed Neumann datum.  This gives their
$C^{2,\alpha}$ bounds.  More generally, if the value relation is
$v_o=r v_i+h$ with $r>0$, and the flux relation after reflection is
$a\partial_tv_i+b\partial_tv_o^{\rm refl}=k$, then the harmonic
combinations $v_o^{\rm refl}-rv_i$ and $av_i+bv_o^{\rm refl}$
have Dirichlet and Neumann data.  Their coefficient determinant is
$a+br>0$.  It stays positive if $b$ decreases to zero with $a>0$.
Variable fixed $C^{2,\alpha}$ ratios are treated by freezing, as in the
next paragraph.  Localization differentiates the value row at most
twice and the flux row at most once; the constants depend only on
the corresponding coefficient norms.  The ordinary boundary and interior estimates follow
from the scalar estimates, with the lapse estimated before the
height row.

For fixed regular variable coefficients, localization gives the same
estimate with lower-order terms.  In a sufficiently small chart
the leading coefficients differ little in supremum norm from their
frozen values.  Apply the constant-coefficient estimate to cutoff
fields on that chart.  In a product such as
$(a-a_*)D^2v$, the coefficient supremum multiplies the top H\"older
seminorm and can be absorbed.  Its fixed H\"older seminorm multiplies
the second-derivative supremum, which is bounded by an arbitrarily
small multiple of the full high norm plus a fixed multiple of the
field supremum, by interpolation.  The same argument applies to
the order-$1$ height coupling and to the differentiated value
relation.  Cutoff commutators have strictly lower orders.  A finite
cover and interpolation for separated-point H\"older quotients give
the claimed local or global variable-coefficient estimate.  This
argument requires only fixed coefficient bounds and is uniform
along compact families of such coefficients.
\end{proof}

We give the nonlinear application of this estimate in detail.  Fix
$0<\alpha'<\alpha<\gamma$, decreasing $\gamma$ in
Lemma~\ref{lem:seam-holder} if necessary.  On a fixed truncation use
the Banach domain
\begin{equation}
 X=C^{3,\alpha'}(\overline\Omega_i)_f
   \times C^{2,\alpha'}(\overline\Omega_i)_{N_i,w_i}
   \times C^{2,\alpha'}(\overline\Omega_{o,R})_{N_o,w_o},
 \label{eq:transmission-domain-space}
\end{equation}
with the two height Dirichlet values imposed by an affine lift of
$u_*$, and with the open condition $N_i,N_o,w_i,w_o>0$.
There are no other conditions built into this domain.  The output
space, denoted $\mathscr Y$, contains the height equation in
$C^{1,\alpha'}$, the two lapse and two conformal-factor equations
in $C^{0,\alpha'}$, the seam value residuals and the far Dirichlet
residuals in $C^{2,\alpha'}$, and the flux residuals at $S$ and
$\mathcal C$ in $C^{1,\alpha'}$.  Height boundary values are not
additional outputs because they are prescribed in $X$.

The quotient in the trace cutoff causes no additional loss.  In a
distance collar $r_S\geq0$ of $S$, with $F=u_*M-f=0$ on $S$,
\begin{equation}
 \frac{F(z',r_S)}{r_S}
   =\int_0^1\partial_{r_S}F(z',\theta r_S)\,\dd\theta.
 \label{eq:transmission-quotient}
\end{equation}
It is consequently an affine bounded map from the height domain
into $C^{2,\alpha'}$.  The same formula controls it in all the
order-by-order estimates below.  Outside this collar its denominator
is smooth and positive.  Products and the finite smooth cutoffs
therefore define a smooth residual map
\begin{equation}
 \mathscr F:[0,1]\times X_+\longrightarrow\mathscr Y,
 \label{eq:transmission-nonlinear-map}
\end{equation}
where $X_+$ denotes the positive open domain after the height lift.

At an exact zero, $w\geq1$.  Indeed its equation has nonnegative
right-hand side, its outer value is $1$, and its natural fluxes
match with zero inner flux.  Testing with $(1-w)^+$ gives the claim.
Since $\Sigma_+x$ is bounded on its support by the finite state
range and the cutoff $x<1-\delta/2$, and $P$ has finite strength,
the forcing in this scalar joined equation is bounded independently
of the particular zero.  Testing with $w-1$ and using the
Dirichlet Poincar\'e inequality gives an $H^1$ bound for $w-1$ on
the fixed truncation.  Local scalar boundedness and H\"older estimates
for its measurable uniformly elliptic joined divergence matrix give
a $C^\gamma$ bound, after reducing $\gamma>0$.  Reflection handles
the homogeneous inner flux.  Thus
\begin{equation}
 \mathscr U=(\dd f,N_i,N_o,w_i,w_o)
 \quad\hbox{is bounded piecewise in }C^\gamma.
 \label{eq:transmission-state-holder-all}
\end{equation}

We now prove an a priori $C^{2,\alpha}$ bound for $\mathscr U$.
For a smooth zero, differentiate the nondivergence height row once.
Its highest terms are the height and lapse terms in
\eqref{eq:transmission-principal-rows}.  Expanding the lapse and
$w$ equations shows that their only highest terms at these column
orders are their own second derivatives.  Every other term involves
at most two factors of first derivatives of $\mathscr U$, with
coefficients that are smooth functions of its bounded state and of
the quotient \eqref{eq:transmission-quotient}.  At the seam,
differentiating the value relation twice tangentially gives
\begin{equation}
 \partial_a\partial_b(N_o-s_i)
  =\partial_a\partial_b N_o
    -D s_i\,\partial_a\partial_b(N_i,\dd f)
    -D^2s_i\bigl(\partial_a(N_i,\dd f),
                 \partial_b(N_i,\dd f)\bigr)
    +\mathcal R_{ab},
 \label{eq:transmission-value-differentiation}
\end{equation}
where $\mathcal R_{ab}$ contains only bounded geometry coefficients
and lower-order derivatives.  The leading first differential is
exactly the principal value relation already used.  One derivative
of a current-flux relation has only the indicated first-order
principal boundary operators; its other terms are lower at these
column orders.  Formula~\eqref{eq:transmission-quotient} gives the
same derivative count for every occurrence of $F/d_S$.

Here are the interpolation exponents needed for absorption, to make
the dependence on the prior continuity estimate explicit.  Put
$H=1+\|\mathscr U\|_{C^{2,\alpha}}$ and suppose
$\|\mathscr U\|_{C^\gamma}\leq C_0$.  On fixed slightly enlarged
patches, H\"older interpolation gives
\begin{align}
 \|D\mathscr U\,D\mathscr U\|_{C^{0,\alpha}}
   &\leq C H^{\theta_1},
 &\theta_1&=\frac{2+\alpha-2\gamma}{2+\alpha-\gamma}<1,
 \label{eq:transmission-product-interpolation}\\
 \|D^2\mathscr U\|_{C^0}
   &\leq C H^{\theta_2},
 &\theta_2&=\frac{2-\gamma}{2+\alpha-\gamma}<1.
 \label{eq:transmission-sup-interpolation}
\end{align}
Terms of still lower order have smaller exponents or admit the usual
$\eta H+C_\eta$ bound.  The first estimate follows by assigning
one first derivative its supremum norm and the other its
$C^{0,\alpha}$ norm; their interpolation numerators are
$1-\gamma$ and $1+\alpha-\gamma$, respectively.

Freeze the coefficients in a chart of radius chosen using only
\eqref{eq:transmission-state-holder-all}.  Their oscillations in
supremum norm are then as small as desired, while their
$C^\gamma$ norms are bounded.  A principal coefficient error
times a top derivative has $C^{0,\alpha}$ norm bounded by a small
constant times $H$ plus a constant times $H^{\theta_2}$.
The same estimate after two tangential derivatives applies to the
value row in \eqref{eq:transmission-value-differentiation}; its
quadratic part is controlled by
\eqref{eq:transmission-product-interpolation}.  Apply
Lemma~\ref{lem:transmission-schauder} to cutoff fields with matched
cutoff values at a seam.  A second, slightly larger cutoff extends
coefficient differences without losing their small supremum norm.
Cutoff commutators contain only lower derivatives and have fixed
patch-dependent bounds.  Taking the maximum over a finite cover,
and using supremum interpolation for H\"older quotients with
separated points, gives
\begin{equation}
 H\leq \eta H+C\bigl(1+H^{\theta_1}+H^{\theta_2}
                         +H^{\theta_3}\bigr),\qquad
 0\leq\theta_3<1.
 \label{eq:transmission-global-absorption}
\end{equation}
Choose $\eta<1/2$.  This bounds $H$.  Constants for differentiated
cutoffs and source functions may depend on the finite parameter
choices, which is permitted here.

We must also justify using smooth-zero estimates for zeros initially
in \eqref{eq:transmission-domain-space}.  The following difference
quotient argument supplies that justification.  Off the seam,
ordinary scalar bootstrapping applies, estimating lapse and $w$
first and height next; at $S$ the $w$ conormal condition is its
ordinary homogeneous Neumann condition because the height gradient
is normal.  Near a seam flatten the face and take a tangential
difference quotient of the equations.  For the value relation use
the averaged first differential of $s$ between the two states.
For a fixed zero these coefficients are uniformly
$C^{2,\alpha'}$ and, on a sufficiently small patch and for
sufficiently small differences, are supremum-close to the same
frozen principal coefficients.  The difference quotients have
uniformly bounded one-lower norms from the starting regularity.
Apply the fixed-coefficient estimate to their cutoffs.  Principal
oscillations times top seminorms are absorbed; top supremum terms
are interpolated against the known one-lower norm.  All
commutators and differentiated explicit data involve only these
known orders.  This controls one tangential derivative in the same
column spaces.

The normal derivatives are recovered in a definite order.  Solve
the lapse and $w$ equations for their double normal derivatives;
their strictly positive normal coefficients permit division.
Then solve the height row for its double normal derivative, using
the newly controlled lapse derivatives.  Other second derivative
terms have at least one tangential index, which is already
controlled.  Iterating proves smoothness.  More explicitly, if
orders $m+3,\alpha'$ for height and $m+2,\alpha'$ for the other
fields are known, apply the same estimate to a tangential
difference quotient of $m$ tangential derivatives.  In the height
row, including its one residual derivative, every distributed
coefficient term other than the principal row on the increments
uses at most $m+2,\alpha'$ derivatives of $\mathscr U$.
The other interior rows start with at most one derivative of
$\mathscr U$ outside their own principal terms.  In the value row
the only new-order term is its first state differential on the
increment.  Thus the same argument gives the next derivatives with
the largest number of tangential indices.  Recover the others with
successively fewer tangential indices by the lapse and $w$
equations first, then the height equation.  These identities hold
on the open side and extend continuously to the face.  The quotient
near $S$ obeys \eqref{eq:transmission-quotient} at every step.

Consequently every zero of \eqref{eq:transmission-nonlinear-map}
is smooth, and \eqref{eq:transmission-global-absorption} bounds it
at the stronger exponent $\alpha$.  On a fixed truncation the
positive zero set over $0\leq u_*\leq1$ is compact in $X$:
the stronger H\"older bounds give precompactness, the residual map
is continuous, and the uniform positive floors for $N$ and $w$
keep limits in $X_+$.

\subsection{Fredholm index and continuation}

\begin{lemma}[Index zero for the linearized transmission problem]
\label{lem:transmission-fredholm}
The field derivative $D_X\mathscr F$ at each zero of
\eqref{eq:transmission-nonlinear-map} is Fredholm of index zero
between \eqref{eq:transmission-domain-space} and its specified
residual space.  Independently, the scalar density and adjoint
problems with smooth base metrics, fixed $C^{1,\alpha'}$ drift,
a fixed positive $C^{2,\alpha'}$ value ratio, natural flux matching,
inner Neumann and outer Dirichlet data are
Fredholm of index zero.  Their estimates may depend on derivatives
of the fixed ratio.
\end{lemma}

\begin{proof}
Discard compact lower-order terms in the linearization.  This is
legitimate with the column orders used above.  For example, a
variation of the trace target in the expanded lapse row contains
at most two derivatives of the height variation or one of the
other variations, and therefore maps compactly to its
$C^{0,\alpha'}$ residual.  Variations of $F/d_S$ lose one height
derivative by \eqref{eq:transmission-quotient}, which still leaves
this compactness.  In contrast, the first differential of
$N_o-s(N_i,\dd f)$ in the $C^{2,\alpha'}$ seam residual is
principal and is retained.  The principal height coupling to a
lapse derivative in its $C^{1,\alpha'}$ residual is also retained.

In this principal operator scale the background $\dd f$ continuously
to zero and recompute $q,\mathcal P,\beta,s,A$, the height coupling,
and the value-differential coefficients from that scaled background.
This scales coefficient fields, not the variations.  In particular
$ck=q-1$ in \eqref{eq:transmission-frozen-relations} holds all
along the path.  The states remain in a fixed positive ellipticity
range.  The local estimates of
Lemma~\ref{lem:transmission-schauder}, followed by localization,
give along this path
\begin{equation}
 \|v\|_X\leq C\bigl(\|L_\lambda v\|_{\mathscr Y}
                                  +\|v\|_0\bigr),
 \label{eq:transmission-linear-apriori}
\end{equation}
where $\|v\|_0$ is a sum of field supremum norms and $C$ is
uniform in the path parameter $\lambda$.

At zero background gradient the height decouples, and the lapse
and $w$ blocks have scalar value and positive conormal transmission.
Decrease each outer derivative coefficient in its flux relation to
zero, leaving the inner coefficient positive.  The determinant
$a+br$ in the proof of
Lemma~\ref{lem:transmission-schauder} stays positive, so
\eqref{eq:transmission-linear-apriori} remains valid.  At the
endpoint each scalar block is triangular: solve the inner Neumann
problem, then the outer Dirichlet problem with the trace prescribed
from the inner solution.  On each connected inner component the
scalar Neumann problem has equal kernel and compatibility-cokernel
dimensions, both one; the outer Dirichlet and height Dirichlet
problems have index zero.  Dividing by positive scalar principal
factors and discarding connection terms reduces this assertion to
ordinary Poisson theory.  Thus the endpoint index is zero.

We justify index propagation using only
\eqref{eq:transmission-linear-apriori}.  A sufficiently fine fixed
finite grid of field evaluations defines
$\mathfrak p:X\to\R^l$ with
$\|v\|_0\leq\eta\|v\|_X+C_\eta|\mathfrak p v|$.
Choose $\eta$ small in \eqref{eq:transmission-linear-apriori}.
The augmented operators
\begin{equation}
 \mathfrak A_\lambda=(L_\lambda,\mathfrak p):
                   X\longrightarrow\mathscr Y\times\R^l
 \quad\hbox{satisfy}\quad
 \|v\|_X\leq C\|\mathfrak A_\lambda v\|.
 \label{eq:transmission-augmented-bound}
\end{equation}
They are injective with closed range.  At an index-zero Fredholm
parameter their range has codimension $l$.  If parameters
$\lambda_j$ of this kind tend to $\lambda$, and the limit range
had codimension at least $l+1$, choose an $(l+1)$-dimensional
subspace $Z_0$ disjoint from that range.  At each $\lambda_j$
there is a unit vector $z_j\in Z_0\cap\operatorname{ran}
\mathfrak A_{\lambda_j}$.  Its preimages $v_j$ are bounded by
\eqref{eq:transmission-augmented-bound}.  A subsequence of $z_j$
converges in the finite-dimensional space to a unit vector $z$,
and operator-norm continuity gives
$\mathfrak A_\lambda v_j\to z$.  Closedness of the limit range
is a contradiction.  Its codimension is therefore finite.

The kernel of $L_\lambda$ is finite dimensional because
$\mathfrak p$ is injective on it.  Its range is closed and has
finite codimension: add the finite-dimensional factor
$\{0\}\times\R^l$ to the closed finite-codimensional augmented
range, and then quotient by that factor.  This sum is closed, and
the resulting quotient is exactly $\operatorname{ran}L_\lambda$.
Thus $L_\lambda$ is Fredholm also at the limit.  Openness and local
constancy of the Fredholm index, together with this closedness
argument starting from the endpoint, propagate index zero along
the whole path.  Restoring compact lower-order terms does not
change it.

For the independent scalar density problem take the fixed positive
$C^{2,\alpha'}$ ratio and fixed $C^{1,\alpha'}$ drift as coefficients.
Its local
principal boundary estimate is the scalar positive-ratio estimate
already proved, and its lower-order drift is compact at the scalar
column order two.  Localization uses only the stated coefficient
norms, so no approximation in a H\"older norm is needed.
The same outer-flux deformation and augmentation
argument applies without reference to a nonlinear zero or to an
a priori estimate for such zeros.  For its adjoint the value ratio
and positive derivative weight exchange their roles; the same
scalar combinations apply.  This proves both independent index
assertions.
\end{proof}

We now complete existence on a fixed truncation.  At $u_*=0$ the
height is identically zero.  Indeed its Dirichlet values are zero,
its trace target is zero, and testing its divergence equation by
$f$ gives $\int p\cdot\dd f=0$.  All scaled sources vanish and
$s_i=N_i$, so the lapse solves the scalar harmonic transmission
problem with ordinary value and normal-derivative matching, inner
flux $-b_*$ in the $n$ direction, and far value $1$.  The weak
form on the joined manifold is coercive after subtracting the
outer boundary value.  Equivalently its index is zero by
Lemma~\ref{lem:transmission-fredholm} and its homogeneous kernel
is zero by energy testing.  It therefore has a unique solution.
Comparison with $1$, using the outward inner flux $+b_*$, gives
$N\geq1$.  Scalar transmission regularity makes it piecewise
smooth.  The remaining equation is
$\divg(N^2\dd w)=0$ with matching value and flux, inner zero
flux, and far value $1$; hence $w=1$.

This initial zero is nondegenerate.  Its homogeneous linearization
is triangular: first
$\divg(N^2\dd\dot f)=0$ with zero height traces, then the
ordinary harmonic lapse transmission problem with homogeneous
boundary data, then the homogeneous $w$ divergence problem with
zero far value.  Energy testing makes each variation zero.
Terms differentiating the source cutoffs carry $u_*$ or vanish
because $\ell=0$; constitutive-coefficient variations in these
rows multiply $\dd f=0$ or $\dd w=0$.  Index zero now implies
invertibility of the full field derivative.

We supply the finite-dimensional continuation argument so that no
orientation or uniqueness assumption at the target parameter is
implicit.  Let $\mathcal K$ be the compact set of positive zeros
of \eqref{eq:transmission-nonlinear-map} over $[0,1]$.  By the
Fredholm assertion, a finite-dimensional subspace
$\mathfrak Z\subset\mathscr Y$ can be chosen so that
\begin{equation}
 D_X\mathscr F(X)+\mathfrak Z=\mathscr Y
 \label{eq:transmission-cokernel-span}
\end{equation}
at every point of $\mathcal K$: choose local cokernel complements
and use a finite cover.  Surjectivity persists on a neighborhood
$\mathcal O$ of $\mathcal K$.  If $Q$ is the quotient map to
$\mathscr Y/\mathfrak Z$, then
\begin{equation}
 \mathcal M=\{(u_*,v)\in\mathcal O:
                               Q\mathscr F(u_*,v)=0\}
 \label{eq:transmission-reduction-manifold}
\end{equation}
is a smooth finite-dimensional manifold of dimension
$\dim\mathfrak Z+1$.  The field derivative of $Q\mathscr F$
is onto, so the parameter projection on $\mathcal M$ is a
submersion.  Its only endpoint faces are therefore $u_*=0$ and
$u_*=1$.  Choose a relatively open neighborhood $\mathcal U$ of
$\mathcal K$ in $\mathcal M$ with compact closure contained in
$\mathcal O$ and with relative boundary disjoint from the zero
set.  This is possible in the finite-dimensional locally compact
manifold.  The norm of $\mathscr F$ has a positive minimum on
that relative boundary.

Apply Sard's theorem simultaneously to $\mathscr F|_{\mathcal M}$,
which takes values in $\mathfrak Z$, and its two endpoint faces.
For a sufficiently small common regular value
$z\in\mathfrak Z$, its preimage in $\mathcal U$ is a compact
one-manifold whose boundary lies only on the endpoint faces.
The implicit function theorem at the unique nondegenerate initial
zero gives exactly one point on the $u_*=0$ face.  Compactness of
$\overline{\mathcal U}$ excludes other points on that face for
small $z$: a sequence of such extra points would limit to an
additional zero, or to the initial zero in its uniqueness
neighborhood.  A compact one-manifold has an even number of
boundary points.  Consequently the $u_*=1$ face is nonempty.
Let a sequence of these regular values tend to zero and use
compactness of $\overline{\mathcal U}$ once more.  Its limiting
point is an exact positive zero at $u_*=1$.  This proves the
finite-truncation assertion of
Proposition~\ref{prop:transmission-existence}.

\subsection{Removal of the outer truncation}

Keep every internal parameter fixed and let $R\to\infty$ through
the solutions just obtained.  The lapse, slope, and low-order seam
bounds are locally uniform by
\eqref{eq:compact-bounds} and Lemma~\ref{lem:seam-holder}.  We
explain why the local upper bound for $w$ is also independent of
$R$, since the fixed-truncation energy bound alone does not imply
this assertion.

Choose a fixed sphere $\mathcal S$ outside all source supports.
On its exterior both $N_o$ and $x=N_ow$ are harmonic.  For $x$
this follows directly from the two equations:
\begin{equation}
 N_o\Delta_{g_o}(N_ow)=-P,
 \label{eq:transmission-exterior-x}
\end{equation}
so the conclusion holds wherever $P=0$, in particular beyond
$\mathcal S$.  Both fields have value $1$ at the truncating
boundary.  Positive radial superharmonic barriers on the AF end
of weight $|x'|^{-\beta}$, with any fixed $0<\beta<1$, bound
their deviations from $1$ by their bounded values at
$\mathcal S$ times such a weight.  Their construction uses
$\Delta_{g_o}|x'|^{-\beta}
=\beta(\beta-1)|x'|^{-\beta-2}
+O(|x'|^{-\beta-3})<0$ at large radius; enlarge $\mathcal S$
if necessary.  At a finite outer edge the zero deviation is
consistent with the positive barrier.

Suppose the supremum of $w$ on the fixed region inside
$\mathcal S$ were unbounded, and divide by this supremum $H_j$.
Exterior harmonic comparison for $x$, together with the global
positive floor for $N_o$, bounds the normalized $w$ everywhere
and gives a decay bound $C|x'|^{-\beta}$ for any subsequential
limit.  On compact sets its scalar forcing divided by $H_j$
tends to zero: $\Sigma_+x$ is bounded by its cutoff and the
finite state bounds, and $P$ has fixed finite strength.  The
coefficient matrices converge locally along a subsequence by the
H\"older estimates for $N$ and $\dd f$.  Scalar local
boundedness, H\"older compactness and weak energy estimates
then give a nonnegative limit $w_\infty$, with a nonzero
maximum on the fixed compact region, satisfying
\begin{equation}
 -\divg(A_\infty\dd w_\infty)=0,
 \qquad [w_\infty]_{\mathcal C}=0,
 \qquad [(A_\infty\dd w_\infty)_n]_{\mathcal C}=0,
 \qquad (A_\infty\dd w_\infty)_n|_S=0.
 \label{eq:transmission-normalized-w-limit}
\end{equation}
It tends to zero at infinity.  This is impossible: on arbitrarily
large compact joined domains, the weak maximum principle with
natural inner flux bounds it by the supremum on the distant
Dirichlet sphere, and that supremum tends to zero.  This argument
uses only the scalar joined divergence form and remains valid for
the limiting piecewise coefficients.  The contradiction proves
the required compact upper bound for $w$.

Without normalization the same exterior barriers show that every
limit has $N_o,x\to1$, and hence $w\to1$.  They also give
derivative estimates of every fixed order on annuli after scaling,
by ordinary scalar harmonic estimates for the fixed smooth AF
metric.  On the compact region containing the sources, the
H\"older estimates and the absorption argument above are uniform
in $R$: take an auxiliary fixed surrounding annulus, whose
derivatives are already bounded by these harmonic estimates, and
use the finite cover of compact patches inside it.  A diagonal
subsequence therefore converges in every fixed smooth norm on
each closed side on compact sets.  Its positive floors and its
boundary and transmission conditions persist.  This proves the
outer-limit assertion and completes
Proposition~\ref{prop:transmission-existence}.
